\subsection{Introduction to Relations} 
\begin{definition}[Relation]
A \term{relation} is a set of ordered pairs. If the pair $(x,y)$ is an element of a relation $R$, we say $x$ is \term{related to} $y$, and write $x\rel{R}y$.
\end{definition}
\begin{Note}Informally, we think of relations as expressing a \makeblank between numbers. The notation $x\rel{R}y$ is meant to parallel things like $x<y$.\end{Note}
\begin{Example}
Let $R=\Set{(2,4), (-3,5), (-4,0), (1,-1), (0,-2)}$.
\begin{enumerate}\setabc
\item Fill in each blank to make a true statement.\lbr%
$\makeblanksmall\rel{R}\makeblanksmall$%
\qquad%
$\makeblanksmall\rel{R}\makeblanksmall$%
\qquad%
$\makeblanksmall\rel{R}\makeblanksmall$%
\item For each statement below, determine whether it is true or false.\lbr
$3\rel{R}5$ \makeblanksmall%
\qquad%
$-3\rel{R}5$ \makeblanksmall%
\qquad%
$-2\rel{R}0$ \makeblanksmall
\end{enumerate}
\end{Example}
\begin{definition}[Domain and Range]
The \term{domain} of a relation is the set of all first entries in the pairs. The \term{range} of a relation is the set of all second entries.

For a relation $R$, $\dom R$ denotes the domain of $R$, and $\rng R$ denotes its range.
\end{definition}
\begin{Example}
Let $R=\Set{(2,4), (-3,5), (-4,0), (1,-1), (0,-2)}$.
\begin{itemize}
\item $\dom R=\Set{\makeblank}$
\item $\rng R=\Set{\makeblank}$
\end{itemize}
\end{Example}
\begin{Example}
Let $S=\Set{(0,1), (0,2), (0,3), (1,2), (1,3), (2,3), (3,3)}$
\begin{itemize}
\item $\dom S=\Set{\makeblank}$
\item $\rng S=\Set{\makeblank}$
\end{itemize}
\end{Example}
\begin{Note}When an element of a set is repeated, we only list it \makeblank[1in]. The range and domain may have \makeblank[1in] elements than the relation, but never \makeblank[1in].\end{Note}